Задача коммивояжёра (Travelling Salesman Problem, TSP) --- одна из классических задач комбинаторной оптимизации. Она требует найти кратчайший возможный маршрут, проходящий через все заданные города ровно один раз и возвращающийся в исходную точку. TSP служит базовой моделью для множества практических задач маршрутизации, логистики и планирования.

В TSP города рассматриваются как вершины графа, а расстояния между ними — как рёбра с весами. Необходимо найти гамильтонов цикл минимальной суммарной стоимости. Поскольку число возможных маршрутов растёт факториально с количеством городов, задача принадлежит классу NP-трудных: известные точные алгоритмы имеют экспоненциальную сложность.

Способы нахождения приближенного решения TSP позволяют получить маршрут, приближенный к верному решению задачи, но имеют меньшую сложность и завершаются за разумное время. Для малых графов применяют динамическое программирование, для больших --- эвристические методы: жадные алгоритмы, генетические алгоритмы, имитация отжига.

Цель данной курсовой работы заключается в изучении приближённого алгоритма решения задачи коммивояжёра на основе минимального остовного дерева и разработке программы для визуализации работы алгоритма.